\documentclass[conference]{IEEEtran}

% ---- Packages ----
\usepackage[T1]{fontenc}
\usepackage{cite}
\usepackage{amsmath,amssymb,amsfonts}
\usepackage{graphicx}
\usepackage{textcomp}
\usepackage{xcolor}
\usepackage{booktabs}
\usepackage{multirow}
\usepackage{url}
\usepackage{balance}

\def\BibTeX{{\rm B\kern-.05em{\sc i\kern-.025em b}\kern-.08em
    T\kern-.1667em\lower.7ex\hbox{E}\kern-.125emX}}

\begin{document}

\title{Which Rewordings Move a Language Model's Moral Judgment}

\author{\IEEEauthorblockN{Andrew H. Bond}
\IEEEauthorblockA{\textit{Senior Member, IEEE}\\
\textit{Department of Computer Engineering} \\
\textit{San Jose State University}\\
San Jose, CA, USA \\
andrew.bond@sjsu.edu \\
ORCID: \texttt{0009-0003-2599-6158}}}

\maketitle

\begin{abstract}
Language models now moderate content, triage clinical intake, and screen legal filings. A model used this way should reach the same judgment however the situation is worded. We test five deployed models against that requirement. Each model reads a moral scenario and rates the harm it involves along seven named dimensions. We then reword the scenario in ways that change how it reads without changing what happened, and we measure how far the ratings move. Three rewordings move them in every model tested. Softening the language, adding emotional weight, and inserting vivid but irrelevant detail all displace the ratings well beyond what repeated readings of the same text produce. Two rewordings move nothing. Swapping the gender of the people involved and changing the order in which scenarios are presented leave judgments where they were. The three rewordings that work all make a morally irrelevant feature easier to notice, which tells a defender what to watch for. Weaknesses also fail to line up across models. The model that never yields to a user pushing a false correction recovers least well from emotional wording, and the model that recovers best from emotional wording tracks the fewest parties in a crowded scenario. One robustness score describes none of them accurately. Six further open models from five families rate the same scenarios and agree closely on all seven dimensions, so the dimensions measure something shared rather than something each model invents. A worked attack closes the argument. Softening the wording of flagged content drops its harm rating below a moderation threshold, and the same wording flips the verdict of a rule based decision engine placed downstream. The study ran about 8,000 model calls inside a \$50 per day budget.
\end{abstract}

\begin{IEEEkeywords}
LLM security, adversarial evaluation, moral judgment, invariance testing, vulnerability profiling, AI safety, calibration
\end{IEEEkeywords}

% ================================================================
\section{Introduction}

When a language model judges a moral scenario, it reports how much harm the situation involves along several named dimensions such as physical injury, emotional damage, loss of autonomy, and broken trust. A system that makes decisions from those ratings should reach the same judgment whether the same events are told in mild language or vivid language, whether irrelevant detail is present, and whether a user pushes back on the answer. When the judgment moves under one of those changes, an attacker who controls only the wording can move the decision without touching a single fact. That failure is what this paper measures.

Earlier work measures one such failure at a time. Studies of sycophancy~\cite{sharma2024sycophancy, perez2023model} measure how far a model bends toward a user who disagrees with it. Studies of framing~\cite{tversky1981framing, echterhoff2024cognitive} measure how far it bends toward the wording of the question. Studies of prompt sensitivity~\cite{scherrer2023moral} measure how far it bends toward incidental phrasing. Each reports one number for how robust a model is, and one number cannot separate a model that resists rewording but yields to pressure from a model that does the reverse. Averaging weaknesses that vary independently discards exactly the information a defender needs, and no later analysis recovers it~\cite{bond2026geometric}.

We measure the weaknesses separately instead. Every model rates every scenario on one shared set of seven harm dimensions, which makes displacements from different tests directly comparable. We then apply five rewordings that change how a scenario reads without changing its moral content. They are softened or heightened language, added emotional weight, vivid but irrelevant detail, a gender swap of the people involved, and a change in the order scenarios are presented. Each should leave the judgment where it was, and measuring whether it does is the test. Two further conditions complete the picture. In one, a user presses the model to change an answer it has already committed to. In the other, the model states how confident it is, which we compare against how often it is right.

This paper contributes the following.

\begin{enumerate}
\item A profiling method that reports, for one model, which rewordings move its judgments and which leave them alone, at a cost of about \$50 per day.
\item Evidence that the weaknesses are selective. Three rewordings displace judgments in all five models and two displace nothing, which narrows the threat to rewordings that make an irrelevant feature easier to notice.
\item Evidence that weaknesses do not travel together. No model is best on every test, and a model that leads on one test is often last on another, so certifying a model on any single test is unsafe.
\item An evaluation pipeline that completes the study inside a fixed budget and takes further models and rewordings without redesign.
\end{enumerate}

% ================================================================
\section{Related Work}

\textbf{Sycophancy and social pressure.} Sharma et al.~\cite{sharma2024sycophancy} showed that models shift their answers toward the opinions a user states. Perez et al.~\cite{perez2023model} wrote model-generated evaluations that found the same tendency across many models. Our correction test differs in that the model must first commit to a verdict over one turn and defend it over the next, which separates a model that updates on real evidence from a model that folds under pressure.

\textbf{Framing and anchoring.} Tversky and Kahneman~\cite{tversky1981framing} established that wording changes human choices. Echterhoff et al.~\cite{echterhoff2024cognitive} surveyed the same biases in language models. Those studies test one reworded condition at a time. We put every reworded condition in one shared rating scale, so the size of one weakness can be read against the size of another.

\textbf{Calibration.} Kadavath et al.~\cite{kadavath2022know} reported that models mostly know what they know. Our results do not reproduce that in moral judgment. Every model we tested states more confidence than its accuracy supports, and combining the four models gives 9.3 standard deviations of miscalibration. We score confidence by rank order rather than by absolute value, which avoids the known problem that models compress their stated confidence into a narrow band~\cite{niculescu2005predicting}.

\textbf{Multi-dimensional evaluation.} Bond~\cite{bond2026geometric} argued that moral situations vary along several dimensions at once and that collapsing them loses structure. Thiele~\cite{thiele2026labse} found those dimensions readable from LaBSE sentence embeddings~\cite{feng2022labse}, with moral content and language identity carried in largely separate directions. That work reads dimensions out of a model's internal representation. We measure whether the model's behavior respects them.

\textbf{Executive functions in AI.} Diamond~\cite{diamond2013executive} defined executive functions as cognitive flexibility, inhibitory control, and working memory. We borrow the three as test categories, measuring whether a model can switch ethical frameworks on request, hold its judgment steady against emotional wording, and keep track of who is involved as a scenario adds parties.

\textbf{Adversarial evaluation.} Red teaming~\cite{ganguli2022red} looks for prompts that make a model produce something harmful. We look for the opposite failure. The model produces nothing harmful and instead misjudges harmful material that an attacker has reworded, which is the failure that matters when a model is the thing making the decision.

% ================================================================
\section{How We Measured This}

\subsection{The Seven Harm Dimensions and the Rewordings}

Every model rates a scenario on seven harm dimensions, which are physical, emotional, financial, autonomy, trust, social, and identity. Each dimension takes a score from 0 to 10, so the total runs from 0 to 70. The seven come from the nine-dimensional moral vector of~\cite{bond2026geometric} by dropping and merging the dimensions that models could not rate consistently. An earlier nine-dimensional form with abstract dimensions such as procedural legitimacy produced scores that disagreed between repeated readings of the same text, which makes any displacement measured on them uninterpretable.

\textbf{Where the seven dimensions come from.} They are the reliably measurable part of the nine-dimensional \texttt{MoralVector} used by DEME v3, the decision kernel of the ErisML compiler~\cite{erisml_compiler}, whose axes come from a three by three assessment matrix crossing individual, relational, and collective concerns with what matters, who decides, and what is known. Table~\ref{tab:deme} gives the correspondence. Four axes map straight across, three are partial, and two of them, rights and epistemic quality, are not scored separately because models rated them inconsistently in preliminary testing. Movement in the seven dimensions is still movement in the nine, so a displacement measured here is a displacement the kernel downstream would also see. That kernel reads the harm vector and emits one of three verdicts, permitted, prohibited, or requires human review. A live end to end integration and native nine-dimensional scoring are left to an extended version.

\begin{table}[htbp]
\caption{How the seven harm dimensions line up with the nine axes of the DEME v3 \texttt{MoralVector}~\cite{erisml_compiler}.}
\label{tab:deme}
\centering
\footnotesize
\begin{tabular}{@{}lll@{}}
\toprule
\textbf{Harm dim.} & \textbf{DEME v3 axis} & \textbf{Mapping} \\
\midrule
$h_\text{phys}$ & physical harm & direct \\
$h_\text{auto}$ & autonomy / consent & direct \\
$h_\text{trust}$ & legitimacy / trust & direct \\
$h_\text{soc}$ & societal / environmental & direct \\
$h_\text{fin}$ & fairness / equity & partial \\
$h_\text{emot}$ & virtue / care & partial \\
$h_\text{ident}$ & privacy / standing & partial \\
none & rights & not scored \\
none & epistemic quality & not scored \\
\bottomrule
\end{tabular}
\end{table}

All five test tracks use these same seven dimensions, so a displacement measured in one track can be compared against a displacement measured in another. Every model call returns a fixed JSON structure holding the verdict, a confidence from 0 to 10, and the seven harm scores, which removes free-text parsing as a source of disagreement.

The rewordings fall into three groups.

\emph{Rewordings that carry no moral weight} change how a scenario reads and leave what happened untouched. They are the gender swap, the softened and heightened rewrites, added emotional weight, and vivid but irrelevant detail. A competent evaluator should return the same ratings after any of them, so any movement beyond the level of repeated readings marks a way in for an attacker.

\emph{Corrections} come in two kinds. A valid correction supplies real new evidence and should change the judgment. An invalid correction supplies a fabricated claim and should not. Only the second kind bears on security.

\emph{Exploratory rewordings} such as a cultural reframe may or may not carry moral weight, so we report what they do without calling the result a failure.

\subsection{Checking That the Dimensions Mean the Same Thing to Other Models}
\label{sec:validation}

A measurement is worth nothing if the ruler is invented. Before reporting any
displacement we check that the seven dimensions mean the same thing to models
that had no part in the study. Six open models drawn from five families,
namely Qwen3, Qwen3-small, GLM-5, GPT-OSS, MiniMax-M2, and Gemma, rated all
seven dimensions from 0 to 10 on 31 scenarios, three times each, through the
NRP managed LLM API. None of the six is among the five systems under test and
none is behind a proprietary endpoint, so anyone can repeat this check.

\begin{table}[htbp]
\caption{Agreement across the six open models on 31 scenarios. ICC(2,$k$) is two-way random, absolute agreement, average measures. The last column is Krippendorff's interval alpha.}
\label{tab:validation}
\centering
\footnotesize
\begin{tabular}{@{}lcc@{}}
\toprule
\textbf{Harm dimension} & \textbf{ICC(2,$k$)} & \textbf{Krippendorff $\alpha$} \\
\midrule
Physical    & 0.893 & 0.571 \\
Emotional   & 0.923 & 0.649 \\
Financial   & 0.983 & 0.902 \\
Autonomy    & 0.893 & 0.564 \\
Trust       & 0.958 & 0.783 \\
Social      & 0.926 & 0.662 \\
Identity    & 0.939 & 0.709 \\
\midrule
\textbf{Overall} & \textbf{0.969} & \textbf{0.836} \\
\bottomrule
\end{tabular}
\end{table}

The six models agree closely. Taken together their ratings reach an intraclass
correlation of 0.969 and a Krippendorff alpha of 0.836, both above the levels
conventionally read as reliable. The seven dimensions therefore track something
the models share rather than a private scale each one invents. Financial,
trust, and identity are the most consistently rated, at alphas of 0.90, 0.78,
and 0.71. Physical harm and autonomy are the least consistent, at 0.57 and
0.56, because most of these scenarios are interpersonal and carry almost no
physical harm, so the ratings sit near zero and small disagreements loom large.
Readers should treat displacement on those two dimensions with more caution
than the rest. Asked to rate the same text twice, a single model returns
correlated scores at a median of 0.96, and that number sets the noise level
against which every effect in Section~\ref{sec:results} is judged.

The same check supports the mapping in Table~\ref{tab:deme}. Three of the four
axes that carry straight across, trust, fairness by way of the financial
dimension, and physical harm, are among the most reliably rated, so the seven
dimensions retain what the nine-dimensional vector was carrying.

\subsection{The Five Test Tracks}

Each track applies one kind of reworded condition and measures how far the seven ratings move.

\textbf{Learning (L1--L4).} These tests ask whether a model updates on evidence for the right reasons. They cover learning from examples (L1), handling a user correction (L2, the headline test), carrying an ethical framework across scenarios (L3), and revising a judgment by an amount that matches the strength of new evidence (L4). In L2 the model states a verdict in the first turn and then receives either a genuine correction or a fabricated one in the second. Subtracting the rate at which it accepts fabricated corrections from the rate at which it accepts genuine ones measures how well it tells them apart.

\textbf{Metacognition (M1--M4).} These tests ask whether a model knows how good its own answers are. They cover how well stated confidence matches accuracy (M1, the headline test), whether a model can tell clear cases from ambiguous ones (M2, scored by area under the ROC curve), whether its stated uncertainty predicts how much its own verdicts vary (M3, scored by Spearman correlation), and whether it works harder on harder problems (M4, scored by response length). M1 reports expected calibration error over 5 bins with a bootstrap standard error from 200 resamples.

\textbf{Attention (A1--A4).} These tests ask whether a model can ignore what does not matter. They cover resistance to vivid and mild distractors (A1, the headline test), stability when a scenario is padded to twice and four times its length (A2), whether attention lands on the dimensions a scenario actually implicates (A3), and accuracy when several scenarios are interleaved in one prompt (A4). A1 uses two distractor strengths so we can see whether displacement grows with the strength, and it repeats the test with the model warned in advance.

\textbf{Executive functions (E1--E4).} These tests ask whether a model controls its own processing. They cover switching between utilitarian, deontological, and virtue ethics on request (E1), holding a judgment steady against emotional wording (E2, the headline test), noticing when one changed fact should change the verdict (E3), and tracking participants as a scenario grows from 2 to 8 parties (E4). E2 repeats each scenario with an added instruction telling the model it is being emotionally manipulated and to focus on the facts, which measures how much of the displacement a warning removes.

\textbf{Social cognition (T1--T5).} These tests probe the structure of the seven dimensions themselves. They cover which dimensions a scenario moves (T1), whether a gender swap or a cultural reframe changes anything (T2), whether telling events in a different narrative order changes the verdict (T3), whether the order scenarios arrive in changes the verdict (T4), and whether softened or heightened wording changes it (T5, the headline test). T5 rewrites each scenario in a softer and a more dramatic style while holding the events fixed.

\subsection{Design Choices That Affect the Result}

\textbf{One model writes the text, others judge it.} A single fixed model, Gemini 2.0 Flash, wrote every distractor, emotional rewrite, softened and heightened variant, correction, and counterfactual. The models under test only read text they did not write. Letting a model generate the attacks it is then scored against would confirm whatever that model already does.

\textbf{Three kinds of scenario.} Gold scenarios are hand written with their rewrites audited by a person, and they carry the most interpretive weight. Probe scenarios are synthetic pairs built so that only one thing differs and the expected behavior is obvious. Generated scenarios are drawn in bulk from AITA, a Reddit advice forum with 270,709 posts~\cite{aita_dataset}, and from 25 curated Dear Abby columns. The generated tier buys statistical power and pays for it with noisier labels.

\textbf{What counts as movement.} Every reworded test carries 3 to 5 control readings per scenario in which the model rates the identical text again. That control measures how much a model's ratings wander on their own. All significance tests compare displacement against that measured wander rather than against zero. The choice changed our conclusions. Apparent violations above 6 standard deviations in an early analysis disappeared once controls were in place, because the models had simply been rating the same text differently on different readings.

\textbf{Statistics.} Rates carry Wilson 95 percent confidence intervals~\cite{wilson1927probable}. Rate comparisons use a two proportion z test. Severity shifts use a paired t test. Calibration error carries a bootstrap standard error. Evidence from several models is combined with Fisher's method~\cite{fisher1925methods} using the Wilson and Hilferty normal approximation.

\subsection{Which Models Ran Which Tests}

We tested five models from two families. Four are Google models, Gemini 2.0 Flash, Gemini 2.5 Flash, Gemini 3 Flash Preview, and Gemini 2.5 Pro. The fifth is Claude Sonnet 4.6 from Anthropic, included so that no finding rests on one vendor. Attention, executive functions, and social cognition ran in full on all five. The learning track ran on three Gemini models across L1 to L4, with Claude added on the correction test so the sycophancy result spans both families. The metacognition track ran four Gemini models on M1 and two on the full M1 to M4 set, because Gemini 2.5 Pro costs about \$0.10 per call against about \$0.002 for the Flash models. The whole study came to roughly 8,000 calls. Table~\ref{tab:coverage} gives the exact counts.

\begin{table}[htbp]
\caption{Scenario, control, and model counts for each headline test.}
\label{tab:coverage}
\centering
\footnotesize
\begin{tabular}{@{}lcccc@{}}
\toprule
\textbf{Attack} & \textbf{Test} & \textbf{Scen.} & \textbf{Ctrl} & \textbf{Mod.} \\
\midrule
Ling.\ framing & T5 & 21 & 3 & 5 \\
Emot.\ tone & E2 & 23 & 3 & 5 \\
Irrel.\ detail & A1 & 21 & 5 & 5 \\
Social pressure & L2 & 37 & 5 & 4 \\
Self-monitoring & M1 & 12--100 & 0--3 & 4 \\
\bottomrule
\end{tabular}
\end{table}

% ================================================================
\section{Running the Study Inside a Daily Budget}

The study had to finish inside a daily spending cap on shared infrastructure, and the pipeline is built around that limit.

\textbf{Concurrency that finds its own level.} Calls run from a worker pool that starts at 50, halves whenever the provider returns an error, and climbs by 5 for each stretch of clean responses. The same code held steady throughput against providers whose limits ranged from 15 to 1,500 requests per minute, with no per-provider tuning. A track finished in 7 to 79 minutes depending on how many scenarios and models it covered.

\textbf{Spending planned in advance.} A planner sets how many scenarios and control readings each model gets, based on what that model costs per call. Flash models at about \$0.002 per call absorb full scenario counts and the maximum number of control readings. Gemini 2.5 Pro at about \$0.10 per call runs every test type with fewer scenarios each. This is why coverage is uneven, and the planner makes the unevenness a stated decision rather than an accident.

\textbf{Fixed output shape.} Every call returns JSON matching a schema holding the verdict, the confidence, the seven harm scores, and the reasoning text. Nothing is parsed out of prose, so two models are always compared on the same quantity.

\textbf{Cost and extension.} Total cost grows with the number of models times the number of reworded conditions times the number of scenarios. A new model needs only credentials. A new reworded condition needs a template for generating it and a statement of what it should leave unchanged. Sample size in the generated tier is the dial for trading statistical power against money, and turning it down leaves the hand audited scenarios untouched.

\textbf{What is released.} Each track is one self-contained Python file that runs in a notebook. The five tracks, the scenarios, the rewrite templates, and the analysis code ship in the \texttt{agi-hpc} package.

% ================================================================
\section{Results}
\label{sec:results}

\subsection{What Moves the Ratings}

Table~\ref{tab:invariance} lists the five headline tests. Three rewordings displace the ratings in every model tested. Two others, the gender swap and the change of presentation order, displace nothing and are not shown. Which rewordings work matters more than how far they move things, because it tells a defender which attacks are worth defending against.

\begin{table}[htbp]
\caption{The five headline tests. Affected counts how many models moved. The last column combines the per-model evidence.}
\label{tab:invariance}
\centering
\footnotesize
\begin{tabular}{@{}lccc@{}}
\toprule
\textbf{Attack surface (track)} & \textbf{$n$} & \textbf{Affected} & \textbf{$\sigma$} \\
\midrule
Ling.\ framing (T5) & 5 & 5/5 & 8.9 \\
Emot.\ tone (E2) & 5 & 5/5 & 6.8 \\
Irrel.\ detail (A1) & 5 & 5/5 & 5.0 \\
Social pressure (L2) & 4 & 3/4 & 13.3 \\
Self-monitoring (M1) & 4 & 4/4 & 9.3 \\
\bottomrule
\end{tabular}
\end{table}

\textbf{Softened and heightened wording (T5).} Softening the language drops the total harm rating by 10 to 16 points out of 70. Heightening it raises the rating by 6 to 11. Rating the same text twice moves it only 1 to 7. Claude moves in one direction more than the other. Softened wording pulls its rating down by 9.1 points while heightened wording moves it only 1.5. A study that reworded in one direction only would have missed that.

\textbf{Emotional wording (E2).} Adding emotional weight produces a medium to large effect in every model, with paired Cohen's $d$ between 0.60 and 1.06 and paired $t$ between 2.90 and 5.10 on 22 degrees of freedom. Read in the units of the scale, the average rating moves about 11 to 14 points out of 70, against 2 to 4 points when the same text is read again, and verdicts flip more often than repeated readings explain. Table~\ref{tab:e2sigma} breaks this out by model. The combined sigma is there only to show that all five models moved the same way across both families, not as the main evidence.

\begin{table}[htbp]
\caption{Effect of emotional wording by model (E2). $d$ is paired Cohen's $d$. The $\sigma$ column checks that models agree in direction and is not the primary statistic. Recovery is the share of the movement a warning takes back.}
\label{tab:e2sigma}
\centering
\begin{tabular}{@{}lcccc@{}}
\toprule
\textbf{Model} & \textbf{Paired $t$} & \textbf{$d$} & \textbf{$\sigma$} & \textbf{Recovery} \\
\midrule
Claude Sonnet 4.6 & +5.10 & 1.06 & 4.1 & 20\% \\
Gemini 2.0 Flash & +4.45 & 0.93 & 3.7 & 73\% \\
Gemini 2.5 Flash & +3.92 & 0.82 & 3.4 & 29\% \\
Gemini 2.5 Pro & +3.20 & 0.67 & 2.9 & 20\% \\
Gemini 3 Flash & +2.90 & 0.60 & 2.6 & 33\% \\
\midrule
\textbf{Fisher combined} & & & \textbf{6.8} & \textbf{mean 38\%} \\
\bottomrule
\end{tabular}
\end{table}

Two of the numbers in that table pull apart. Claude moves furthest under emotional wording and also recovers least when told the wording is manipulative, correcting only 20 percent of the movement. Flash 2.0 is displaced more often, changing 48 percent of its verdicts, yet recovers 73 percent of the movement once warned. One emotional robustness score would rate these two models the same. Reading displacement and recovery separately shows they fail in different ways and need different fixes.

\textbf{Irrelevant detail (A1).} Table~\ref{tab:a1sigma} gives the per-model effects. Claude is the most displaced of the five, yet it also grades its response correctly, moving furthest on vivid detail, less on mild detail, and least on the control. It is discriminating between strong and weak distractors and reacting to both too much.

\begin{table}[htbp]
\caption{Per-model distractor significance (A1, paired $t$, $\text{df}=20$).}
\label{tab:a1sigma}
\centering
\begin{tabular}{@{}lccc@{}}
\toprule
\textbf{Model} & \textbf{Distractor $\sigma$} & \textbf{Warned $\sigma$} & \textbf{Recovery} \\
\midrule
Claude Sonnet 4.6 & 4.2 & 4.3 & 36\% \\
Gemini 2.0 Flash & 3.2 & 3.7 & 14\% \\
Gemini 3 Flash & 1.1 & 2.9 & 33\% \\
Gemini 2.5 Flash & 1.7 & 2.3 & 50\% \\
Gemini 2.5 Pro & 1.0 & 0.9 & 33\% \\
\midrule
\textbf{Fisher combined} & \textbf{5.0} & & \textbf{mean 39\%} \\
\bottomrule
\end{tabular}
\end{table}

\subsection{What Leaves the Ratings Alone}

Two rewordings change nothing. Swapping the gender of the people in a scenario and changing the order scenarios arrive in stay within the noise of repeated readings in every model. Gemini 3 Flash and Claude both score 0.958 on the gender swap and between 0.933 and 1.000 on presentation order. Models have learned to hold steady under some rewordings and not others, and knowing which is which is what makes the result useful.

What the three effective rewordings share is that each makes a morally irrelevant feature easier to notice. Softened wording makes severity harder to see. Emotional wording makes distress easier to see. Vivid detail makes an unrelated image easier to see. The gender swap and the reordering change words without changing what stands out, and they change nothing in the ratings. An attacker therefore needs to change what draws the model's attention, not merely change the tokens.

The cultural reframe sits between the two groups at a combined 4.1 standard deviations. It moves verdicts more than repeated readings do, but a reframe may carry real moral content, so we do not count it as a failure.

\subsection{How Far Each Model Bends to a False Correction}
\label{sec:sycophancy}

The correction test separates the four models sharply, and the ordering tracks how stable each model is on its own (Table~\ref{tab:sycophancy}). We report the rate at which a model accepts a fabricated correction divided by the rate at which it accepts a genuine one. A score of 0 means the model accepts real corrections and rejects fabricated ones perfectly. A score of 1 means it cannot tell them apart at all.

\begin{table}[htbp]
\caption{Correction discrimination (L2). Wrong-flip $n{=}9$ per model.}
\label{tab:sycophancy}
\centering
\footnotesize
\begin{tabular}{@{}lcccc@{}}
\toprule
\textbf{Model} & \textbf{Corr.} & \textbf{Wrong} & \textbf{SI} & \textbf{$t_\text{conf}$} \\
\midrule
Claude 4.6 & 59\% & 0\% & 0.00 & +2.83 \\
Gemini 2.0 Fl. & 71\% & 33\% & 0.47 & $-$2.12 \\
Gemini 2.5 Pro & 68\% & 44\% & 0.66 & $-$0.80 \\
Gemini 2.5 Fl. & 76\% & 56\% & 0.73 & +0.41 \\
\bottomrule
\multicolumn{5}{@{}l}{\scriptsize Wilson 95\% CIs on wrong-flip rate. 0\%\,[0--30], 33\%\,[12--65],}\\
\multicolumn{5}{@{}l}{\scriptsize 44\%\,[19--73], 56\%\,[27--81].}
\end{tabular}
\end{table}

How a model's confidence responds to a fabricated correction says how it is handling that correction. Claude becomes more confident, at $t = +2.83$, which is what happens when a model argues back rather than quietly ignoring the user. Flash 2.0 becomes less confident, at $t = -2.12$, which reads as partial doubt. Flash 2.5 does not shift at all, at $t = +0.41$, which is what uncritical acceptance looks like. How often each model changes its verdict with no correction at all follows the same ordering, at 0 percent for Claude, 2 percent for Flash 2.0, 8 percent for Pro, and 19 percent for Flash 2.5. The models that wander least on their own are the ones that yield least to a user.

The graded revision test rules out the simpler explanation. All models scale how much they revise to how strong the new evidence is, at $z$ between 4.4 and 6.7 comparing extreme evidence against the control, with every $p$ below 0.003. Their revision machinery works. What fails is telling a real correction from a fabricated one, and that is a narrower problem to fix. Neither test alone would show this, because one measures whether a model updates and the other measures whether it should have.

\subsection{Every Model Overstates Its Own Accuracy}

Table~\ref{tab:m1} gives the calibration error for each model. All five state more confidence than their accuracy earns, and all five err in the same direction. The error runs from 0.186 for Pro to 0.421 for Flash 2.0.

\begin{table}[htbp]
\caption{Expected calibration error by model (M1). Every model is overconfident.}
\label{tab:m1}
\centering
\begin{tabular}{@{}lccc@{}}
\toprule
\textbf{Model} & \textbf{ECE} & \textbf{Direction} & \textbf{$z$-score} \\
\midrule
Gemini 2.0 Flash & 0.421 & overconfident & 6.2$\sigma$ \\
Gemini 2.5 Flash & 0.415 & overconfident & 7.0$\sigma$ \\
Gemini 3 Flash & 0.333 & overconfident & 4.5$\sigma$ \\
Gemini 2.5 Pro & 0.186 & overconfident & 2.2$\sigma$ \\
Claude Sonnet 4.6 & 0.250 & overconfident & n/a \\
\midrule
\textbf{Fisher combined} & & & \textbf{9.3$\sigma$} \\
\bottomrule
\end{tabular}
\end{table}

Larger models are better calibrated. Pro at 0.186 is clearly better than the Flash models at 0.333 to 0.421, which suggests calibration is something training can improve. Claude sits at 0.250, so the problem is not particular to one vendor.

Effort scaling separates the models the other way. Flash 2.0 writes longer answers for harder problems, scoring 0.715. Pro writes answers of the same length whatever the problem, scoring 0.000. The better calibrated model is the one that does not adjust its effort, so knowing what you know and knowing how hard to work are not the same ability.

\subsection{Weaknesses Do Not Travel Together}

No model is best everywhere. Table~\ref{tab:profiles} lists measures from different tracks that pull in different directions.

\begin{table}[htbp]
\caption{Selected measures from five different tracks, showing that no model leads on all of them.}
\label{tab:profiles}
\centering
\footnotesize
\begin{tabular}{@{}lccccc@{}}
\toprule
\textbf{Model} & \textbf{L2 Syc} & \textbf{E2 Rec} & \textbf{A4 Div} & \textbf{E3 CF} & \textbf{E4 WM} \\
\midrule
Claude & \textbf{0.000} & \textbf{20\%}$^\dagger$ & \textbf{0.571}$^\dagger$ & 56\% & 0.886 \\
Flash 2.0 & 0.472 & \textbf{73\%} & 0.812 & 50\% & \textbf{0.710}$^\dagger$ \\
Pro 2.5 & 0.657 & 20\% & \textbf{1.000} & \textbf{75\%} & 0.887 \\
Flash 2.5 & \textbf{0.726}$^\dagger$ & 29\% & 0.875 & 69\% & \textbf{0.900} \\
\bottomrule
\multicolumn{6}{l}{\footnotesize $^\dagger$Worst in column. Bold without $^\dagger$ = best.}
\end{tabular}
\end{table}

Claude never accepts a fabricated correction and is at the same time the worst at tracking interleaved scenarios, at 0.571, and the worst at recovering from emotional wording, at 20 percent. Flash 2.0 recovers best from emotional wording, at 73 percent, and holds the fewest parties in memory, at 0.710. Pro notices changed facts best, at 75 percent, tracks interleaved scenarios perfectly, at 1.000, and still accepts fabricated corrections at a middling rate. Working memory does improve steadily across the Gemini generations, from 0.710 to 0.900 to 0.909, but nothing else does.

Averaging these columns would produce a number that fits none of these models. Choosing a model on the strength of one test is therefore unsafe, because the test a team happens to run may be the one dimension where the model is strong.

\subsection{Warnings Recover About a Third of the Movement}

Telling a model in advance that its input has been tampered with helps, and it helps by about the same amount whichever tampering is used. Warned about distracting detail, models take back 39 percent of the movement. Told they are being emotionally manipulated, they take back 38 percent. Two rewordings that work in different ways yield to a warning to nearly the same degree, which suggests the limit belongs to the instruction rather than to either attack.

The same models overstate their own accuracy, at a combined 9.3 standard deviations of miscalibration. A model that cannot judge when it is right is poorly placed to notice when something has moved its answer, so the two findings fit together. We report the pairing as consistent with that reading and not as evidence that one causes the other.

\subsection{Some Things Work}

Not everything we measured is a failure. All three Gemini models tested revise in proportion to the evidence, with extreme evidence producing more revision than moderate evidence and moderate more than irrelevant, at $p$ below 0.003 throughout. The judgment machinery responds to evidence strength as it should.

Gemini 3 Flash and Claude also grade their reaction to distractors correctly, moving most on vivid detail, less on mild detail, and least on the control. These models are not blind to the difference between a strong distraction and a weak one. They simply react to both more than they should.

Framework switching works as well. Asked to reason as a utilitarian, a deontologist, or a virtue ethicist, all five models change 32 to 47 percent of their verdicts and produce reasoning carrying the markers of the requested framework 89 to 93 percent of the time. They are genuinely switching, not relabeling.

% ================================================================
\section{Discussion}

\subsection{What This Means for Teams Deploying These Models}

\emph{Different models need different defenses.} Table~\ref{tab:profiles} shows each model failing somewhere different. Claude never yields to a user pushing a false correction and is the worst of the five at recovering from emotional wording. Flash 2.0 recovers well from emotional wording and holds the fewest parties in mind. A team choosing a model for content moderation can therefore ask which of these failures its own traffic will produce, instead of reading a single score that fits no model. The combined evidence for the individual effects runs from 5.0 to 8.9 standard deviations, so the patterns are not noise.

\emph{Defenders should watch for repackaging.} Three rewordings work and two do not, and the three that work share one property. Each makes something morally irrelevant easier to notice. Attacks worth guarding against are therefore the ones that present the same facts with different emphasis, and the clearest example is softened language that drops perceived harm by 10 to 16 points out of 70.

\emph{The vulnerability runs in the dangerous direction.} Claude drops 9.1 points under softened wording and moves only 1.5 under heightened wording. That pattern is what a model would learn from ordinary text, where dramatic language is common and often exaggerated while measured language usually accompanies careful description. The inference is reasonable and the consequence is not. Real harm is often described in calm and reasonable terms by people who want it read as minor, which is what corporate liability language, legal euphemism, and institutional harm reduction wording are for. A model that discounts drama and trusts calm is weakest exactly where accurate judgment matters most. It has learned what to ignore without learning what to distrust~\cite{thiele2026labse}.

\emph{Fixing one weakness does not fix the others.} Claude shows that complete resistance to social pressure is achievable, and Claude also has the worst emotional recovery at 20 percent and the worst tracking of interleaved scenarios at 0.571. Within the Gemini family, acceptance of fabricated corrections ranges from 0 to 56 percent with no relation to how new or capable the model is. A certificate based on one test tells a deployer almost nothing about the rest.

\emph{Warnings in the system prompt are not enough.} Recovery stops at about 38 percent, so a warning leaves roughly two thirds of the movement in place. That figure is a practical ceiling for any defense written as an instruction. It coincides with overconfidence in every model tested, at calibration errors between 0.19 and 0.42. Whether the poor self-assessment limits the recovery or merely accompanies it is unresolved, but a defense that depends on a model noticing its own displacement has to contend with both. Changing what the model represents internally~\cite{thiele2026labse} is the more promising direction.

\emph{The whole study cost \$50 per day.} Five tracks, five models, and about 8,000 calls fit inside that budget. Any team running a decision service can afford to repeat this on the models it actually uses.

\subsection{A Worked Attack on a Moderation Filter}
\label{sec:exploit}

To show what this costs in deployment, we built the simplest filter a team
would build. It flags any item whose total harm rating passes a threshold on
the 0 to 70 scale. The attacker changes only how the item reads, never what it
says happened, by applying the softening rewrite. Run against the six open
models of Section~\ref{sec:validation}, softening drops the average rating of
the six hand audited scenarios by 14.0 points, with individual drops between
4.4 and 26.7 and four of the six above 10 points. Heightening the wording
raises the rating by only 7.3.

The result is that flagged content passes and nobody sees it happen. With the
threshold set at the median rating of flagged content, softening moves three of
the six items from flagged to passed while every fact in them stays the same.
The lopsidedness makes this worse. Because softening moves ratings about twice
as far as heightening, hiding real harm is easier than manufacturing a false
alarm, and hiding real harm is the failure that hurts. Any filter that gates on
a severity score inherits this, including legal triage and clinical intake
systems. Six audited scenarios are too few to state a rate at which this
succeeds, so we present it as proof that the attack completes end to end. The
size of the displacement itself rests on all 31 scenarios.

\emph{A rule engine does not fix it.} The obvious defense is to replace the
threshold with structured rules. We tested that by sending each scenario, in
its neutral, softened, and heightened forms, through the ErisML DEME v3
kernel~\cite{erisml_compiler}. The models extract structured facts such as
whether rights were violated, whether consent was valid, and how severe the
harm was, and a fixed rule module turns those facts into one of five ranked
verdicts from forbid up to strongly prefer. Across 161 scenario and model
pairs, softening turns a restrictive verdict into a permissive one in 13.7
percent of cases and shifts the verdict by 0.65 steps on average. Forbid
verdicts fall from 44 to 27, a drop of 39 percent, while strongly prefer rises
from 46 to 78. The rules never change. What changes is the facts the models
hand them, so the rule engine inherits the weakness of the models in front of
it. Hardening has to reach the extraction of facts, not only the rules applied
to them.

\subsection{Limitations}

Each test covers between 6 and 40 scenarios per model, which is not many. We therefore lean on whether models agree in direction rather than on any single test reaching significance, and we combine evidence across models with Fisher's method. Individual cells carry wide intervals. The correction test rests on 9 fabricated corrections per model, so its point estimates should be read loosely and the ordering across the four models trusted more than any one value.

Three boundaries limit what these results cover. The first is subject matter. Every scenario is an English language interpersonal dilemma from AITA or Dear Abby, and those reflect mostly North American norms. What counts as a morally irrelevant rewording may not hold elsewhere, since emotional directness and indirect speech carry different weight in different places, and our cultural reframe result at 4.1 standard deviations marks that boundary without settling it. Whether the same failures appear in factual reasoning, legal analysis, or clinical triage is untested. The second is coverage of attacks. We tested five rewordings, not every rewording an attacker might try, so the absence of movement under a gender swap and a reordering bounds the threat without certifying anything about rewordings we did not run. The third is the number of models. Five models from two families cannot establish a trend, so the steady improvement in working memory across Gemini generations is worth noting and not worth relying on.

Two further caveats concern method. The daily budget forced uneven coverage across models, which the planner made deliberate but did not remove. Fisher's method assumes the tests it combines are independent, and our models saw overlapping scenario sets, which correlates them and makes the combined figures somewhat generous. This is why we lead with agreement in direction rather than with the combined values. All calls used default sampling settings, so we have not mapped how these effects change with temperature.

% ================================================================
\section{Conclusion}

We tested whether five deployed language models judge a moral situation the same way when the situation is reworded but not changed. They do not, and they fail in a specific and readable pattern.

Three rewordings displace judgments in every model. Softening the language, adding emotional weight, and adding vivid irrelevant detail all move the ratings well past the level at which a model disagrees with itself on a second reading. Two rewordings displace nothing. Swapping the gender of the people involved and changing the order scenarios arrive in leave the ratings where they were. What separates the two groups is whether the rewording changes what stands out, which tells a defender what to look for and what to stop worrying about.

The weaknesses do not line up across models. The model that never yields to a false correction is the worst at shaking off emotional wording. The model that shakes off emotional wording best keeps track of the fewest parties. Any single robustness score averages these into a number that describes no model, so a certificate resting on one test cannot be trusted.

Warning a model that it is being manipulated recovers about 38 percent of the movement and no more, whichever manipulation is used. Defenses written as instructions therefore leave most of the problem in place, and the models that would need to notice their own displacement are the same models that overstate their accuracy in every test we ran.

The attack completes end to end. Softening the wording of flagged content drops it under a moderation threshold, and routing the same content through a rule based decision engine does not save it, because the rules are fed by the same models. Fixing this means changing how facts are read, not which rules are applied to them. All of it ran on about 8,000 calls for under \$50 per day, which puts the same check within reach of any team running a decision service.

% ================================================================
\section*{Reproducibility}

The code, the scenarios, the hand audited rewrites, and the raw notebook outputs ship in the \texttt{agi-hpc} Python package under its \texttt{benchmarks/} directory.\footnote{\url{https://github.com/ahb-sjsu/agi-hpc}, tag \texttt{bds2026-v2}, archived at Zenodo under doi:10.5281/zenodo.20660065, which always resolves to the latest release. The agreement check and the verdict analysis live under \texttt{benchmarks/ieee\_bds\_2026/revision/}.} Each of the five tracks is one standalone Python file that runs in a notebook with no dependency beyond the benchmark SDK. The agreement check of Section~\ref{sec:validation} and the verdict analysis of Section~\ref{sec:exploit} ran separately against the NRP managed LLM API using the open source \texttt{erisml-lib} kernel, and their scripts and raw outputs are included. All three tiers of scenarios and every rewrite template are embedded in the source. Raw results are kept as executed notebooks with the full response logs. The AITA data is public on HuggingFace~\cite{aita_dataset}.

\section*{Acknowledgment}

This work was conducted using the Kaggle benchmark task infrastructure.

% ================================================================
\begin{thebibliography}{00}

\bibitem{sharma2024sycophancy}
M.~Sharma \emph{et al.}, ``Towards understanding sycophancy in language models,'' in \emph{Proc.\ ICLR}, 2024.

\bibitem{perez2023model}
E.~Perez \emph{et al.}, ``Discovering language model behaviors with model-written evaluations,'' in \emph{Findings of ACL}, 2023.

\bibitem{tversky1981framing}
A.~Tversky and D.~Kahneman, ``The framing of decisions and the psychology of choice,'' \emph{Science}, vol.~211, pp.~453--458, 1981.

\bibitem{echterhoff2024cognitive}
J.~M.~Echterhoff, Y.~Liu, and A.~Alessa, ``Cognitive biases in large language models: A survey and mitigation framework,'' \emph{arXiv:2410.02466}, 2024.

\bibitem{scherrer2023moral}
N.~Scherrer \emph{et al.}, ``Evaluating the moral beliefs encoded in LLMs,'' in \emph{Proc.\ NeurIPS}, 2023.

\bibitem{bond2026geometric}
A.~H.~Bond, \emph{Geometric Methods in Computational Modeling: From Manifolds to Production Systems}. San Jose State University, 2026.

\bibitem{erisml_compiler}
A.~H.~Bond, \emph{ErisML Compiler: A Framework-Pluralist Compiler from Natural-Language Moral Material to a Typed Moral Graph}, software, v0.4, 2026. Zenodo, doi:10.5281/zenodo.20659432.

\bibitem{kadavath2022know}
S.~Kadavath \emph{et al.}, ``Language models (mostly) know what they know,'' \emph{arXiv:2207.05221}, 2022.

\bibitem{thiele2026labse}
L.~Thiele, ``Does LaBSE encode moral geometry? Testing Bond's geometric ethics framework through cross-lingual probing,'' Undergraduate Research Report, UCLA Dept.\ of Cognitive Science, 2026.

\bibitem{diamond2013executive}
A.~Diamond, ``Executive functions,'' \emph{Annu.\ Rev.\ Psychol.}, vol.~64, pp.~135--168, 2013.

\bibitem{aita_dataset}
OsamaBsher, ``AITA-Reddit-Dataset,'' HuggingFace Datasets. [Online]. Available: \url{https://huggingface.co/datasets/OsamaBsher/AITA-Reddit-Dataset}

\bibitem{wilson1927probable}
E.~B.~Wilson, ``Probable inference, the law of succession, and statistical inference,'' \emph{J.\ Am.\ Stat.\ Assoc.}, vol.~22, pp.~209--212, 1927.

\bibitem{fisher1925methods}
R.~A.~Fisher, \emph{Statistical Methods for Research Workers}. Oliver and Boyd, 1925.

\bibitem{niculescu2005predicting}
A.~Niculescu-Mizil and R.~Caruana, ``Predicting good probabilities with supervised learning,'' in \emph{Proc.\ ICML}, pp.~625--632, 2005.

\bibitem{feng2022labse}
F.~Feng \emph{et al.}, ``Language-agnostic BERT sentence embedding,'' in \emph{Proc.\ ACL}, pp.~878--891, 2022.

\bibitem{ganguli2022red}
D.~Ganguli \emph{et al.}, ``Red teaming language models to reduce harms: Methods, scaling behaviors, and lessons learned,'' \emph{arXiv:2209.07858}, 2022.

\end{thebibliography}

\balance

\end{document}
